%% ---------------------------------------------------------------------------
%% proposal.tex
%%
%% Research Proposal, main document.
%%
%% ---------------------------------------------------------------------------
\documentclass[12pt,letterpaper]{article}
\usepackage[english]{babel}     % supports english, but default is
% \usepackage[spanish]{babel}
% include this if you want to import graphics files with /includegraphics

\usepackage{longtable}
\usepackage{ifpdf}
\usepackage[table]{xcolor}

\usepackage{anysize}
\marginsize{2.5cm}{2.5cm}{1cm}{1cm}
\usepackage{textcomp}
\usepackage{url}
\bibliographystyle{unsrt}
\usepackage{graphics}
\usepackage{amssymb}
\usepackage{graphicx}
%\usepackage{slashbox}
\usepackage[latin1]{inputenc}
\usepackage{tikz}
\usepackage{amsmath}
\usepackage{float}
\usetikzlibrary{arrows,positioning}
% Color and strikethrough



\usepackage{color}
\usepackage{soul}

\usepackage{array}
\usepackage{makecell}

\usepackage{sectsty}
\allsectionsfont{\sffamily}

\definecolor{dblue}{RGB}{0,102,153}
\newcommand{\dB}[1]{\textcolor{dblue}{\textbf{#1}}}



\setlength{\parskip}{1em}

% Nombre del Estudiante
\newcommand{\scriptAuthor}{Daniel Moya S�nchez}
 
% T��tulo de la tesis
\newcommand{\scriptTitle}{Progress report \#6 for the project: Design of (ASIPs) for Approximate Computing} 


% Keywords
\newcommand{\scriptKeywords}{key, words, ...}

% Para el PDF (cambiar si se desea otras cosas a lo indicado arriba
\newcommand{\pdfAuthor}{\scriptAuthor}
\newcommand{\pdfTitle}{\scriptTitle} 
\newcommand{\pdfKeywords}{\scriptKeywords}


\tikzset{
    mynode/.style={rectangle,rounded corners,draw=black, top color=white, bottom color=yellow!50,very thick, inner sep=1em, minimum size=3em, text centered},
    myarrow/.style={->, >=latex', shorten >=1pt, thick},
    mylabel/.style={text width=7em, text centered} 
} 

\begin{document}
 
 \graphicspath{{./}{./fig/}}

 %% ---------------------------------------------------------------------------
%% titlepage.tex
%%
%% Title page
%%
%% ---------------------------------------------------------------------------

\thispagestyle{empty} 

\begin{center}

\textsc{\LARGE Instituto Tecnol\'ogico de Costa Rica} \\
\textsc{\Large \'Area Acad\'emica de Ingenier\'ia en Computadores}

\textsc{\Large Proyecto de Dise\~no en Ingenier\'ia en Computadores}


\par\vspace{20mm}

\includegraphics[scale=0.24]{logoTEC}

\par\vspace*{\fill}

{\LARGE\bf{\textsf{ \Huge \scriptTitle}}}

\par\vspace*{\fill}

Chair for Embedded System (CES) \\
Kalrsruhe Institute of Technology (KIT) \\
Period: 19/02/2018 (week 2) - 23/02/2018 (week 4) 
%Master Thesis {\sf Proposal} \\ 
%in fulfillment of the requirements for the degree of
%Plan de Proyecto
%Master of Science in Electronics Engineering \\
%Emphasis on Embedded Systems

\par\vspace{20mm}

\textsc{\Large \scriptAuthor}

\vspace*{\fill}

{\today}

\end{center}
\newpage 
\cleardoublepage  


%  \tableofcontents

 \clearpage

\section{Performed activities} \label{act}

\begin{enumerate}
  
 \item \textbf{Implement the ASIPs in the error tolerant applications found (ID 03)}: The \emph{sob} assembly instruction was implemented for the Sobel algorithm.
 This allows computing the following operation in a single cycle:
 \begin{align}
  rd = a*a + b*b
 \end{align}
 The operation described in (1) allows the execution of two multiplications and an addition operation, which is used frequently in the Sobel algorithm. Other operations in this algorithm could not be
 turned into special instructions (for example, a matrix multiplication) because they require more than two paremeters. 
 
 
 \item \textbf{Implement two small assembly codes for each special instruction developed (ID 10)}: The manually-written version of the assembly code for the \emph{sob} instruction was implemented with the
 same format as the previous two assembly codes (explained in the previous report).
 
 
 \item \textbf{Compare execution time, area and power vs error in selected applications (ID 09)}: Comparison at a simulation level was performed, with an amount of lower total cycles found for the 
programs with the special instruction implemented (as expected). With the \emph{sob} instruction, nearly a 53\% reduction of the total cycles was achieved. Power and area were also compared, but
little difference was found in the results. The special instructions used around 1\% more area and around 0.005W more power, due to the fact that the special instructions do not add much 
hardware compared to the processor in general. This results were only simulated since there is not a hardware platform available for the implementation. 

 \end{enumerate}
 

\section{Scope Changes} \label{change}

The total number of approximate applications has been re-modified to 3, with the last one being the \emph{sob} instruction implemented during this last period. 


\section{Earned Value analysis}

The project is almost on schedule, the tasks that were worked on are at 90\% because approval from the supervisor is being waited.
In general, the tasks have taken less effort than expected, mainly due to the fact that the scope has been modified because of all the difficulties encountered through the project. 

Table \ref{tab:gain} summarizes the gained value analysis. 

\begin{table}[h!]
\begin{center}
\caption{Earned Value} \bigskip
\resizebox{\textwidth}{!}{\begin{tabular}{|c|c|c|c|c|c|c|c|c|c|c|c|c|c|} 
 \hline
\makecell{Activity\\ID}	&Activity	&Budget	&\makecell{\%Planned\\Value}	&PV	&AC	&\makecell{\%Completed\\work}	&EV	
&CPI	&SPI	&\makecell{Initial planned\\date}	&\makecell{Ending\\date}	&\makecell{Initial real\\date}	&\makecell{Real ending} \\ \hline


01	&\makecell{Get to know \\the software\\ platform}	&32	&100\%	&32	&34	&100\%	&$32$
&$0.94$	&1	&Week 1	&Week 3	&Week 1	&Week 7 \\ \hline

02	&\makecell{Find appropiate\\ error-tolerant\\ Applications}	&32	&100\%	&32	&20	&100\%	&32
&1.6	&1	&Week 4	&Week 6	&Week 4	&Week 10 \\ \hline

03	&\makecell{Implement the\\ ASIPs in the\\ error tolerant\\ applications\\ found}	&64	&100\%	&64	&40	&90\%	&57.6
&1.44	&0.9	&Week 7	&Week 12	&Week 9	&- \\ \hline


04	&\makecell{Write Project\\Plan Document}	&8	&100\%	&8	&10	&100\%	&8
&0,8	&1	&Week 1	&Week 2	&Week 1	&Week 3 \\ \hline


05	& \makecell{Write\\Requirements\\Document}	&8	&100\%	&8	&6	&100\%	&8	
&1,33	&1	&\makecell{Week 2}	&\makecell{Week 3}	&\makecell{Week 2}	&\makecell{Week 4} \\ \hline


06	& \makecell{Write Design\\ Document}	&8	&100\%	&8	&7	&100\%	&8
&1.14	&1	&\makecell{Week 3}	&\makecell{Week 4}	&\makecell{Week 3}	&\makecell{Week 5} \\ \hline

09	& \makecell{Compare execution\\ time, area and\\ power vs error\\ in selected \\applications}	&32	&75\%	&24	&15	&90\%	&28.8
&1.92	&1.2	&\makecell{Week 12}	&\makecell{Week 16}	&\makecell{Week 12}	&\makecell{-} \\ \hline

10	& \makecell{Implement two \\small assembly \\codes for each\\ special instruction\\ developed}	&16	&100\%	&16	&12	&90\%	&14.4
&1.2	&0.9	&\makecell{Week 11}	&\makecell{Week 13}	&\makecell{Week 11}	&\makecell{-} \\ \hline

& Total	& &	&192	&144	&	&188.8
&1.31	&0.98	&	&	&	& \\ \hline
\end{tabular}}
\label{tab:gain}
\end{center}
\end{table}




\section{Difficulties Encountered} \label{dif}

\begin{itemize} 

\item The server nodes in which the tools like ASIPMeister and Dlxsim work properly were turned off daily at 12:00 pm in Germany (4:00 pm in Costa Rica) for the reported period, so the 
time frame for possible work was reduced. 
 
\end{itemize}




\section{Hard Skills Required/Acquired}

\begin{itemize}
 \item Knowledge in the following software frameworks has been acquired: ASIPMeister and Dlxsim. 
 
 \item Knowledge in HDL has been reinforced.
\end{itemize}




\section{Soft Skills Required/Acquired}

The following soft skills have been excercised:

\begin{itemize}
 
 \item Self-Motivation: Given that there is no direct round-the-clock supervision, self-motivation has been key in working 
 continously with the ASIPMeister and Dlxsim on the creation of special instructions for the approximate applications.
 
 \item Communication: Remote communication has been performed with Jorge Castro for the guidance of this project.
 
\end{itemize}


\section{Lessons Learned}
There are no lessons learned for the reported period. 


% Referencias del Background y el Related Work
\bibliographystyle{sty/plainurl}
\bibliography{references}



\end{document}

